\heading{Activity 1: ask, inspect, verify, iterate}{Work with Claude \hfill Save your work at minute 40}
\topic{Questions worth asking Claude}
\begin{itemize}
  \item ``Before changing anything, inspect this folder and tell me what you think should and should not be tracked in Git.''
  \item ``Explain what \code{.gitignore} is and recommend what this project should ignore.''
  \item ``Explain what a commit represents in plain language.''
  \item ``Show me how to check what Git thinks has changed before we commit anything.''
  \item ``I don't understand the last command. Explain it without assuming I know Git.''
  \item ``Before making this change, tell me how I could undo it if something goes wrong.''
\end{itemize}

\topic{You are done when}
\checkitem{I know which folder is the repository root.}\par
\checkitem{Git is set up in that folder, and not inside another repository.}\par
\checkitem{I understand broadly which files Git tracks.}\par
\checkitem{Large or inappropriate files are left out where appropriate.}\par
\checkitem{I have a first commit with the original files as received, before any change.}\par
\checkitem{I can say how my local repository differs from GitHub.}\par
\checkitem{The repository is connected to GitHub, if you are ready.}\par
\checkitem{I can ask Claude to explain the current state of the repository.}

\topic{If something goes wrong}
\begin{itemize}
  \item ``The push failed. Explain the error and my options.''
  \item ``I think I opened the wrong folder. How do I check where the repository root is?''
  \item ``Show me what changed since the last commit, and whether anything unexpected is included.''
  \item If Claude asks permission at every step: ``Do the next three steps, then stop and summarize.''
  \item If it did a lot at once: ``Explain each thing you just did, in order, and how I would undo each one.''
\end{itemize}
Ask for help after five minutes stuck. A local repository with the push pending is an acceptable checkpoint.

\topic{Finished early?}
Stay in the same conversation. Ask Claude to help you run the notebook in a recorded environment, check its results, and explain what a reproducible baseline is. Then choose: extract tested functions, or extend the analysis in one direction (\href{https://github.com/dowlinglab/claude-for-researchers/blob/main/resources/workshops/cstr/activities/workshop1_extensions.md}{extension ideas}).

\topic{Reflect}
What did Claude explain that you would not have known to ask? Where did you accept a step you did not understand? What did you check yourself?

\vfill
\textbf{My next question or step:}\par\vspace{8pt}\hrule\vspace{18pt}\hrule
